\chapter{Abdominal Pain}
\index{Abdominal Pain}

\section*{Definition and Overview}
Abdominal pain is a symptom defined as discomfort or distress felt anywhere between the chest and the groin. It is one of the most frequent complaints in both general practice and emergency medicine, and its causes span an enormous range, from trivial and self-limiting upsets to life-threatening surgical and vascular emergencies. The pain may arise within the abdominal organs themselves, the peritoneum, the pelvis, or the abdominal wall, but it can also be referred from the chest, spine, or elsewhere. Careful attention to the site, character, onset, radiation, and associated features of the pain is central to identifying the underlying cause and to deciding whether the situation is urgent.

\section*{Epidemiology}
Abdominal pain accounts for a large proportion of primary care consultations and a significant share of emergency department presentations worldwide. Many episodes are short-lived and benign: viral gastroenteritis, food-borne illness, and constipation are among the commonest causes. Appendicitis affects approximately 7\% of people during their lifetime, gallstones occur in around 1 in 10 adults, and chronic conditions such as irritable bowel syndrome, functional dyspepsia, and peptic ulcer disease are highly prevalent. Although the majority of cases resolve without intervention, a small but important minority represent surgical, obstetric, or vascular emergencies in which delay in diagnosis substantially worsens outcomes.

\section*{Aetiology and Pathophysiology}
Abdominal pain arises through several mechanisms. Visceral pain originates from the hollow organs and is typically dull, poorly localised, and felt in the midline, often accompanied by autonomic symptoms such as nausea and sweating. Parietal pain results from inflammation of the peritoneum and is sharper, more localised, and aggravated by movement. Referred pain is felt at a site distant from its source because of shared spinal segments, such as shoulder-tip pain in diaphragmatic irritation.

\begin{itemize}
    \item \textbf{Infection and inflammation:} gastroenteritis, appendicitis, diverticulitis, cholecystitis, pancreatitis, gastritis, and pelvic inflammatory disease
    \item \textbf{Obstruction:} bowel obstruction, severe constipation, and stones blocking the biliary or urinary tracts
    \item \textbf{Ulceration and mucosal disease:} peptic ulcer and inflammatory bowel disease, including Crohn's disease and ulcerative colitis
    \item \textbf{Vascular causes:} mesenteric ischaemia and a leaking or ruptured abdominal aortic aneurysm
    \item \textbf{Gynaecological causes:} dysmenorrhoea, ovarian cysts, endometriosis, and ectopic pregnancy
    \item \textbf{Urinary causes:} urinary tract infection and renal colic
    \item \textbf{Functional disorders:} irritable bowel syndrome and functional dyspepsia, where no structural abnormality is identifiable
    \item \textbf{Metabolic causes:} diabetic ketoacidosis, uraemia, and rare conditions such as acute porphyria
    \item \textbf{Referred sources:} pneumonia, myocardial infarction, shingles, and spinal pathology presenting as abdominal pain
\end{itemize}

\section*{Clinical Features}
The clinical picture must be characterised carefully, as the details of the pain guide diagnosis.

\begin{itemize}
    \item \textbf{Site:} upper, central, or lower abdomen, and whether the pain has migrated, as in appendicitis moving to the right lower quadrant
    \item \textbf{Character:} cramping or colicky, burning, tearing, or constant and gnawing
    \item \textbf{Onset:} sudden and severe (suggesting perforation, rupture, or torsion) versus gradual and progressive
    \item \textbf{Radiation:} to the back, flank, shoulder, or groin
    \item \textbf{Aggravating and relieving factors:} relationship to eating, movement, coughing, or urination
    \item \textbf{Associated symptoms:} fever, vomiting, diarrhoea, constipation, rectal bleeding, urinary symptoms, or jaundice
    \item \textbf{Examination findings:} localised tenderness, guarding, rigidity, abdominal distension, and changes in bowel sounds
\end{itemize}

\section*{Differential Diagnosis}
The differential diagnosis of abdominal pain is broad and organised by region and tempo of presentation.

\begin{itemize}
    \item \textbf{Surgical emergencies:} appendicitis, perforated peptic ulcer, bowel obstruction, volvulus, and ruptured abdominal aortic aneurysm
    \item \textbf{Infective causes:} gastroenteritis, diverticulitis, cholecystitis, cholangitis, and urinary tract infection
    \item \textbf{Gynaecological and obstetric:} ectopic pregnancy, ovarian cyst complications, pelvic inflammatory disease, and dysmenorrhoea
    \item \textbf{Gastrointestinal disease:} peptic ulcer, gallstones, pancreatitis, inflammatory bowel disease, and irritable bowel syndrome
    \item \textbf{Referred and extra-abdominal:} pneumonia, myocardial infarction, shingles, and renal or spinal pathology
\end{itemize}

\section*{When to Seek Medical Care}
A doctor should be consulted if pain is persistent, recurrent, or worsening, or if it is accompanied by fever, vomiting, weight loss, change in bowel habit, jaundice, or pain on urination. Abdominal pain in pregnancy, in older adults, or in people with weakened immunity warrants particularly careful assessment.

\textbf{Contact your local emergency services for:}
\begin{itemize}
    \item Sudden, severe, or rapidly worsening abdominal pain
    \item A hard, rigid, or very tender abdomen, especially with vomiting or fever
    \item Vomiting blood, or passing black, tarry, or bloody stools
    \item Abdominal pain accompanied by chest pain, difficulty breathing, or collapse
    \item Abdominal pain in pregnancy, or in a person who appears pale, sweaty, and unwell
\end{itemize}

\section*{Diagnosis and Investigations}
Diagnosis begins with a thorough history and physical examination, and is refined by targeted investigation.

\begin{itemize}
    \item \textbf{History:} detailed description of onset, site, radiation, character, and associated symptoms, including menstrual, drug, and travel history
    \item \textbf{Examination:} abdominal inspection, palpation, percussion, and auscultation, together with pelvic and rectal examination where indicated
    \item \textbf{Blood tests:} full blood count, inflammatory markers, liver function tests, pancreatic enzymes, renal function, and a pregnancy test in women of reproductive age
    \item \textbf{Urine testing:} for infection, blood, or ketones
    \item \textbf{Imaging:} ultrasound, computed tomography, or plain radiography depending on the suspected cause
    \item \textbf{Endoscopy:} gastroscopy or colonoscopy for upper or lower gastrointestinal pathology
    \item \textbf{Laparoscopy:} considered when the diagnosis remains uncertain despite other investigations
\end{itemize}

\section*{Management}
Management depends on the underlying cause, and the priority is always to identify and treat serious conditions without delay.

\begin{itemize}
    \item \textbf{Treat the cause:} antibiotics for infection, surgical removal for appendicitis or gallstones, and acid suppression for peptic ulcer
    \item \textbf{Symptom relief:} simple analgesia, rest, and clear fluids while the cause is being established
    \item \textbf{Surgical treatment:} required for appendicitis, bowel obstruction, perforation, and other emergencies
    \item \textbf{Dietary measures:} light, low-fat meals for gallstone or dyspeptic symptoms; fibre, fluids, and laxatives for constipation
    \item \textbf{Functional disorders:} reassurance, dietary modification, stress management, and occasionally prescription medicines for irritable bowel syndrome
    \item \textbf{Follow-up:} review to confirm resolution and to detect any new or persisting features
\end{itemize}

\section*{Complications}
\begin{itemize}
    \item Delayed or missed diagnosis of serious causes such as appendicitis, perforation, or a leaking aneurysm
    \item Perforation of the bowel leading to peritonitis and sepsis
    \item Chronic pain and impaired quality of life when functional disorders are poorly managed
    \item Recurrent attacks of gallstones, kidney stones, or peptic ulcers
    \item Complications of treatment, including surgery and the side effects of long-term analgesia
\end{itemize}

\section*{Prognosis and Outcomes}
The outlook depends entirely on the underlying cause. Most episodes of abdominal pain, particularly those from gastroenteritis and other benign causes, resolve within days with supportive care. Surgical conditions such as acute appendicitis carry an excellent prognosis when treated promptly. Chronic functional conditions such as irritable bowel syndrome often improve substantially with dietary and lifestyle management, although they may fluctuate over time. Recognising and treating serious causes early is the single most important factor in achieving good outcomes.

\section*{Prevention and Self-Care}
\begin{itemize}
    \item Practise good hand hygiene and safe food preparation to reduce the risk of gastroenteritis
    \item Eat a fibre-rich diet and drink adequate fluids to prevent constipation
    \item Maintain a healthy weight, lowering the risk of gallstones and reflux disease
    \item Limit alcohol, which can irritate the stomach and provoke pancreatitis
    \item Do not ignore persistent or recurring pain; seek medical review early
    \item Manage stress, which can aggravate functional gut symptoms
    \item Keep up to date with recommended screening, such as cervical screening in eligible women
\end{itemize}

\section*{Sources and Further Reading}
The following sources provide reliable information about the causes, assessment, and management of abdominal pain for both clinicians and patients.

\begin{itemize}
    \item NHS. \textit{Abdominal pain: overview, causes, and when to get help.} https://www.nhs.uk/conditions/abdominal-pain/
    \item Mayo Clinic. \textit{Abdominal pain: symptoms and causes.} https://www.mayoclinic.org/diseases-conditions/abdominal-pain/
    \item MSD Manual. \textit{Abdominal pain: professional evaluation and causes.} https://www.msdmanuals.com/
    \item MedlinePlus. \textit{Abdominal pain health information.} https://medlineplus.gov/abdominalpain.html
    \item NICE. \textit{Gastrointestinal conditions and abdominal pain guidance.} https://www.nice.org.uk/guidance/conditions-and-diseases/
    \item World Health Organization. \textit{Health topics: gastrointestinal and digestive health.} https://www.who.int/
\end{itemize}

\section*{Definition and Overview}
Abdominal pain in children refers to discomfort or distress felt in the belly of an infant, child, or adolescent. It is among the most common reasons families present to general practitioners and emergency departments. The overwhelming majority of causes are benign and self-limiting, but a small number are surgical, infective, or otherwise serious, and the range of likely causes shifts markedly with age. Because young children cannot describe their symptoms accurately, clinicians and carers must rely on careful observation of how the child looks, behaves, feeds, sleeps, and passes urine and stool.

\section*{Epidemiology}
Abdominal pain is extremely common throughout childhood. Functional abdominal pain, in which no organic cause is found, affects roughly 1 in 10 school-aged children and is more prevalent than any single structural cause. Constipation and infant colic are very common in the first years of life, and gastroenteritis causes frequent brief episodes at all ages. Appendicitis is uncommon under the age of 5 and peaks in older children and adolescents. Intussusception, a form of bowel obstruction in which one segment of intestine telescopes into another, occurs most often between 6 and 18 months. Urinary tract infections are an important cause of abdominal pain, particularly in infants and young girls.

\section*{Aetiology and Pathophysiology}
The causes of childhood abdominal pain are best considered by age group, since the anatomy, physiology, and exposure of the gut differ across development.

\begin{itemize}
    \item \textbf{Infants:} infant colic, gastroenteritis, constipation, urinary tract infection, cow's milk protein allergy, and gastro-oesophageal reflux
    \item \textbf{Bowel obstruction and torsion:} intussusception, malrotation with volvulus, and incarcerated hernias can obstruct or strangulate the bowel
    \item \textbf{Young children:} gastroenteritis, mesenteric adenitis, constipation, and urinary infection
    \item \textbf{Older children and adolescents:} appendicitis, inflammatory bowel disease, constipation, and functional abdominal pain
    \item \textbf{Adolescent girls:} dysmenorrhoea, ovarian cyst complications, pelvic inflammatory disease, and rarely ectopic pregnancy
    \item \textbf{Boys:} testicular torsion, which can present with abdominal and lower pelvic pain
    \item \textbf{Functional causes:} visceral hypersensitivity and disordered gut-brain interaction, often triggered or aggravated by stress, anxiety, or previous infection
    \item \textbf{Referred pain:} extra-abdominal sources such as lower-lobe pneumonia may be perceived as abdominal pain
\end{itemize}

\section*{Clinical Features}
Assessment begins with the pattern of pain, but in younger children behaviour is often more informative than words.

\begin{itemize}
    \item Site of pain, and whether the child is able to point to a specific spot
    \item In babies, inconsolable crying, drawing the knees up to the chest, or poor feeding
    \item Vomiting, particularly green (bile-stained) or blood-stained vomit
    \item Blood in the stool, which in intussusception classically resembles redcurrant jelly
    \item Fever, diarrhoea, or constipation
    \item Reduced intake of fluids or food, and reduced urine output
    \item Abdominal distension, tenderness, or a rigid abdomen on examination
    \item Pain or burning with urination, suggesting urinary tract infection
    \item In older children, pain that interrupts sleep, school, or play
\end{itemize}

\section*{Differential Diagnosis}
\begin{itemize}
    \item Gastroenteritis versus appendicitis versus mesenteric adenitis in school-aged children
    \item Intussusception versus gastroenteritis in infants and toddlers
    \item Constipation versus functional abdominal pain in young children
    \item Urinary tract infection versus other abdominal causes at any age
    \item Testicular torsion versus other causes of lower abdominal pain in boys
    \item Gynaecological causes versus appendicitis in adolescent girls
    \item Organic disease versus functional pain or anxiety-related symptoms
\end{itemize}

\section*{When to Seek Medical Care}
Parents should seek urgent review if a child has bile-stained vomiting, blood in the stool, a swollen or tender abdomen, a high fever, inability to keep down fluids, or pain severe enough to interfere with normal activity.

\textbf{Contact your local emergency services for:}
\begin{itemize}
    \item A child who looks seriously unwell: pale, floppy, drowsy, or not responding normally
    \item Difficulty breathing or abnormally rapid breathing
    \item A rigid, tense, or very tender abdomen, or green vomit
    \item Blood in the stool combined with severe crying and drawing up of the legs
    \item Sudden, severe pain with collapse or fainting
\end{itemize}

\section*{Diagnosis and Investigations}
Diagnosis relies on a structured history from the parent or carer, close observation, and targeted tests.

\begin{itemize}
    \item \textbf{History:} onset and site of pain, feeding and drinking, urine output, bowel habit, fever, and any change in behaviour or activity
    \item \textbf{Examination:} vital signs, hydration status, growth, and gentle abdominal palpation, with examination of the groins in boys
    \item \textbf{Urine testing:} dipstick and culture to exclude urinary tract infection
    \item \textbf{Blood tests:} full blood count and inflammatory markers where infection or appendicitis is suspected
    \item \textbf{Ultrasound:} the key investigation for suspected intussusception and helpful in appendicitis and ovarian pathology
    \item \textbf{Plain X-ray or CT:} used where bowel obstruction, perforation, or other serious pathology is possible
    \item \textbf{Family and psychosocial assessment} where functional abdominal pain is suspected
\end{itemize}

\section*{Management}
Management is directed at the underlying cause, with supportive care while the diagnosis is being clarified.

\begin{itemize}
    \item \textbf{Supportive care:} adequate fluids, rest, and paracetamol or ibuprofen at weight-appropriate doses
    \item \textbf{Gastroenteritis:} oral rehydration solution and continuation of normal feeding
    \item \textbf{Constipation:} increased dietary fibre, fluids, and laxatives under medical supervision
    \item \textbf{Appendicitis:} prompt surgical removal of the appendix
    \item \textbf{Intussusception:} pneumatic or hydrostatic reduction under imaging, with surgery reserved for failure or complications
    \item \textbf{Urinary tract infection:} appropriate antibiotics and follow-up where recurrent
    \item \textbf{Functional pain:} reassurance, a regular routine, and psychological support such as cognitive behavioural therapy; medication is rarely required
    \item \textbf{Follow-up:} review to ensure symptoms settle and that growth, sleep, and school attendance return to normal
\end{itemize}

\section*{Complications}
\begin{itemize}
    \item Dehydration from vomiting and reduced fluid intake
    \item Perforation and peritonitis if appendicitis or bowel obstruction is missed
    \item Recurrence of intussusception, particularly after non-surgical reduction
    \item Loss of a testis if torsion is not treated as an emergency
    \item Ongoing functional pain, school absence, and anxiety if the condition is not managed well
    \item Kidney scarring from recurrent or unrecognised urinary tract infections
\end{itemize}

\section*{Prognosis and Outcomes}
Most episodes of abdominal pain in childhood resolve quickly and completely. Appendicitis treated early has an excellent prognosis, and intussusception is usually cured by reduction. Functional abdominal pain frequently improves over months with reassurance and consistent support, although a minority of children continue to experience symptoms into adolescence. The key determinant of outcome is early recognition of serious causes, so prompt medical review of any child who looks unwell is strongly encouraged.

\section*{Prevention and Self-Care}
\begin{itemize}
    \item Encourage thorough hand washing and safe food handling to reduce gastroenteritis
    \item Provide a fibre-rich diet and adequate water to prevent constipation
    \item Maintain regular routines for meals, sleep, and toileting
    \item Keep immunisations up to date, as they protect against several gut infections
    \item Watch for changes in toilet or feeding habits and raise concerns with a doctor
    \item Support children's emotional wellbeing, since stress and anxiety can worsen tummy pain
\end{itemize}

\section*{Sources and Further Reading}
The following trusted organisations provide further reading on childhood abdominal pain for parents, students, and clinicians.

\begin{itemize}
    \item NHS. \textit{Stomach ache in children: causes and when to get help.} https://www.nhs.uk/conditions/stomach-ache/
    \item Mayo Clinic. \textit{Abdominal pain in children: causes and when to see a doctor.} https://www.mayoclinic.org/diseases-conditions/abdominal-pain/
    \item MedlinePlus. \textit{Abdominal pain in children.} https://medlineplus.gov/ency/article/003120.htm
    \item MSD Manual. \textit{Abdominal pain in children: professional reference.} https://www.msdmanuals.com/
    \item Royal Children's Hospital Melbourne. \textit{Clinical practice guidelines on abdominal pain.} https://www.rch.org.au/clinicalguide/
    \item Centers for Disease Control and Prevention. \textit{Childhood health and development resources.} https://www.cdc.gov/
\end{itemize}
